% from ega2-1-fr.tex
\oldpage[II]{181}
et par suite $g''\circ u$ est un morphisme projectif
(\hyperref[II.5.5.5-fr]{5.5.5, (ii)})~; par la commutativité de
(\hyperref[II.8.8.4.1-fr]{8.8.4.1}), il en est de même de $v\circ g$, et comme
$v$ est affine, donc séparé, on en conclut finalement que $g$ est projectif
(\hyperref[II.5.5.5-fr]{5.5.5, (v)}).
\end{proof}

\begin{env}[8.8.5]
\label{II.8.8.5-fr}
Reprenons la situation de (\hyperref[II.8.8.1-fr]{8.8.1}). Nous allons voir
qu'on peut donner du morphisme $g:\mathbf{V}(\mathcal{L})\to C$ une autre
définition valable pour tout $\mathcal{O}_X$-Module inversible $\mathcal{L}$
(non nécessairement ample). Pour cela, considérons le $f$-morphisme
\begin{equation}
\label{II.8.8.5.1-fr}
  \tau^\flat:\mathcal{S}\longrightarrow
  \bigoplus_{n\geqslant 0}\mathcal{L}^{\otimes n}
\tag{8.8.5.1}
\end{equation}
correspondant au morphisme $\tau$ de
(\hyperref[II.8.8.1.2-fr]{8.8.1.2}). On en déduit
(\hyperref[II.1.5.6-fr]{1.5.6}) un morphisme $g':V\to C$ tel que, si
$\pi:V\to X$ et $\psi:C\to Y$ sont les morphismes structuraux, les diagrammes
\begin{equation}
\label{II.8.8.5.2-fr}
  \xymatrix{
    X\ar[d]_{f} & V\ar[l]^{\pi}\ar[d]^{g'}\\
    Y & C\ar[l]_{\psi}
  }
  \qquad
  \xymatrix{
    X\ar[r]^{j}\ar[d]_{f} & V\ar[d]^{g'}\\
    Y\ar[r]_{\varepsilon} & C
  }
\tag{8.8.5.2}
\end{equation}
soient commutatifs
(\hyperref[II.8.5.1.2-fr]{8.5.1.2} et
\hyperref[II.8.5.1.3-fr]{8.5.1.3}). Nous allons montrer que (si l'on suppose
$\mathcal{L}$ ample pour $f$) les morphismes $g$ et $g'$ sont identiques.

En effet, la question est locale sur $Y$, et on peut donc supposer que
$Y=\operatorname{Spec}(A)$ est affine, et (en raison de
(\hyperref[II.8.8.1.3-fr]{8.8.1.3})) identifier $X$ à un ouvert de
$P=\operatorname{Proj}(S)$, où
\[
  S=A\oplus\bigoplus_{n\geqslant 1}
  \Gamma(X,\mathcal{L}^{\otimes n})~;
\]
on déduit alors de (\hyperref[II.8.8.1.4-fr]{8.8.1.4}) que
$\Gamma(X,\mathcal{O}_P(n))=\Gamma(X,\mathcal{L}^{\otimes n})$ pour tout
$n\in\mathbf{Z}$. Tenant compte de la définition de
$h=\operatorname{Spec}(\alpha)$, où $\alpha$ est le $p$-morphisme canonique
$\widetilde{S}\to\mathcal{S}_P^{\geqslant}$
(\hyperref[II.8.6.1.2-fr]{8.6.1.2}), il faut vérifier que la restriction à
$X$ de
$\alpha^\sharp:p^*(\widetilde{S})\to\mathcal{S}_P^{\geqslant}$ est
identique à $\tau$. Compte tenu de
(\hyperref[0.4.4.3-fr]{0, 4.4.3}), on est ramené à voir que si on compose
l'homomorphisme canonique
$\alpha_n:S_n\to\Gamma(P,\mathcal{O}_P(n))$ avec l'homomorphisme de
restriction
\[
  \Gamma(P,\mathcal{O}_P(n))\longrightarrow
  \Gamma(X,\mathcal{O}_P(n))=
  \Gamma(X,\mathcal{L}^{\otimes n}),
\]
on obtient l'identité pour tout $n>0$~; or, cela résulte aussitôt de la
définition de l'algèbre $S$ et de celle de $\alpha_n$
(\hyperref[II.2.6.2-fr]{2.6.2}).
\end{env}

\begin{proposition}[8.8.6]
\label{II.8.8.6-fr}
Supposons (avec les notations de (\hyperref[II.8.8.5-fr]{8.8.5})) que si
l'on pose $f=(f_0,\lambda)$, l'homomorphisme
$\lambda:\mathcal{O}_Y\to f_*(\mathcal{O}_X)$ soit bijectif~; alors~:
\begin{enumerate}
  \item[\rm{(i)}] Si on pose $g=(g_0,\mu)$,
  $\mu:\mathcal{O}_C\to g_*(\mathcal{O}_V)$ est un isomorphisme.
  \item[\rm{(ii)}] Si $X$ est intègre (resp. localement intègre et normal),
  $C$ est intègre (resp. normal).
\end{enumerate}
\end{proposition}

\begin{proof}
En effet, le $f$-morphisme $\tau^\flat$ est alors un \emph{isomorphisme}
\[
  \tau^\flat:\mathcal{S}=\psi_*(\mathcal{O}_C)
  \longrightarrow f_*(\pi_*(\mathcal{O}_V))
  =\psi_*(g_*(\mathcal{O}_V))
\]
et le $Y$-morphisme $g$ peut être considéré comme celui pour lequel
l'homomorphisme $\mathcal{A}(g)$
(\hyperref[II.1.1.2-fr]{1.1.2}) est égal à $\tau^\flat$. Pour voir que
$\mu$ est un isomorphisme de $\mathcal{O}_C$-Modules, il suffit
(\hyperref[II.1.4.2-fr]{1.4.2}) de voir que
\[
  \mathcal{A}(\mu):\psi_*(\mathcal{O}_C)
  \longrightarrow\psi_*(g_*(\mathcal{O}_V))
\]
est un isomorphisme. Mais par définition
(\hyperref[II.1.1.2-fr]{1.1.2}), on a
$\mathcal{A}(\mu)=\mathcal{A}(g)$, d'où la conclusion de (i).

Pour prouver (ii), on peut se ramener au cas où $Y$ est affine, donc
$\mathcal{S}=\widetilde{S}$, avec

